\vspace{-2mm}

\section{结论}
\vspace{-2mm}


本文重新审视了在语言模型蒸馏中同策略采样真正必要的情形。通过系统比较同策略蒸馏（OPD）与提出的半同策略蒸馏（Semi-OPD）在多种教师-学生组合上的表现，我们发现持续重新生成学生模型的轨迹通常并非必需，甚至在教师与学生对齐程度较差时可能造成负面影响。我们的分析揭示了一个关键原则：有效的蒸馏需要在学生和教师两个层面平衡同策略性，初始的教师-学生重叠度可作为判断何时采用OPD或Semi-OPD的实用指标。除了提供一种比OPD更高效的替代方案外，我们的研究还推动未来工作将同策略采样视为一种可选择的设计策略而非默认选项，并开发能够显式考虑教师-学生对齐程度与轨迹质量的蒸馏方法。


